\documentclass[12pt]{article}
\usepackage[T1]{fontenc}
\usepackage[utf8]{inputenc}
\usepackage{lmodern}
\usepackage{textcomp}
\usepackage{enumitem}
\usepackage{geometry}
\geometry{margin=1in}
\begin{document}
\begin{center}
Answer Key - Astronomy - Graduate Level Assessment 6 \
\small (Answers are ordered exactly as in the question paper.)
\end{center}

\section*{Multiple Choice}
\begin{enumerate}[label=\arabic*.]
\item \textbf{Answer:} B
\\ \textit{Options:} A) Green Bank Telescope; B) Arecibo Telescope; C) Yevpatoria RT-70 Telescope; D) Effelsberg Telescope
\\ \textit{Note:} The Arecibo Telescope, with its 305-meter (1,000-foot) diameter parabolic antenna, was the largest and most powerful radio telescope in the world until its collapse in 2020, making it the correct answer for this question.
\item \textbf{Answer:} B
\\ \textit{Options:} A) Being hit in the head by a bullet.; B) Being hit by a small meteorite.; C) Starvation during global winter caused by a major impact.; D) Driving while intoxicated without wearing seatbelts.
\\ \textit{Note:} Being hit by a small meteorite is the least likely cause of death due to the extremely low frequency of such events occurring and the generally small size and mass of small meteorites, which significantly reduces their potential for causing fatal injuries.
\item \textbf{Answer:} B
\\ \textit{Options:} A) The materials they are made of are not normally gaseous in everyday experience.; B) They actually contain a significant fraction of non-gaseous matter.; C) The materials that make up these planets are primarily in the form of a plasma not a gas.; D) Actually it's a great description because these worlds are big and gaseous throughout.
\\ \textit{Note:} The term "gas giants" is misleading because, although Jupiter and the other jovian planets are primarily composed of hydrogen and helium gases, they also contain substantial amounts of non-gaseous materials such as ices, rocks, and metals, which contribute to their overall mass and structure.
\item \textbf{Answer:} D
\\ \textit{Options:} A) hydrogen compounds (H$_{2}$OCH$_{4}$NH$_{3}$) light gases (H He) metals rocks; B) hydrogen compounds (H$_{2}$OCH$_{4}$NH$_{3}$) light gases (H He) rocks metals; C) light gases (H He) hydrogen compounds (H$_{2}$OCH$_{4}$NH$_{3}$) metals rocks; D) light gases (H He) hydrogen compounds (H$_{2}$OCH$_{4}$NH$_{3}$) rocks metals
\\ \textit{Note:} The correct answer lists the ingredients of the solar nebula in order of abundance, with light gases such as hydrogen and helium being the most prevalent, followed by hydrogen compounds, then rocks, and finally metals, reflecting the compositional gradient observed in the early solar system formation.
\item \textbf{Answer:} D
\\ \textit{Options:} A) Altair A*; B) Alsephina A*; C) Fomalhaut A*; D) Sagittarius A*
\\ \textit{Note:} The correct answer is Sagittarius A* because it is the name designated to the supermassive black hole located at the center of the Milky Way galaxy, identified through various astronomical observations and studies.
\end{enumerate}

\section*{True / False}
\begin{enumerate}[label=\arabic*.]
\setcounter{enumi}{5}
\item \textbf{Answer:} False
\\ \textit{Options:} A) True; B) False
\\ \textit{Note:} The Great Red Spot is a massive storm located in Jupiter's atmosphere, not Saturn's, which is characterized by its own distinct features such as rings and storms but does not have the Great Red Spot.
\item \textbf{Answer:} True
\\ \textit{Options:} A) True; B) False
\\ \textit{Note:} The altitude of the north celestial pole corresponds directly to your latitude in the northern hemisphere because, at your location, the angle between the horizon and the north celestial pole is equal to the angle between the horizon and the equator, which is defined by your latitude.
\item \textbf{Answer:} True
\\ \textit{Options:} A) True; B) False
\\ \textit{Note:} The statement is true because the ecliptic represents the plane of Earth's orbit around the Sun and corresponds to the Sun's projected path against the background of stars as observed from Earth throughout the year.
\item \textbf{Answer:} True
\\ \textit{Options:} A) True; B) False
\\ \textit{Note:} Pluto is indeed accompanied by one medium-sized satellite, Charon, which is approximately half the size of Pluto, along with two smaller satellites, Nix and Hydra, which have been confirmed through observations and data from the Hubble Space Telescope.
\item \textbf{Answer:} False
\\ \textit{Options:} A) True; B) False
\\ \textit{Note:} Dark energy is a mysterious form of energy that permeates space and accelerates the expansion of the universe, rather than being related to radiation emitted by black holes or their interactions with matter.
\end{enumerate}

\section*{Long Form Response}
\begin{enumerate}[label=\arabic*.]
\setcounter{enumi}{10}
\item \textbf{Answer:} As one travels north from the United States into Canada, the apparent altitude of Polaris, the North Star, increases. This phenomenon occurs because Polaris is located nearly directly above the North Celestial Pole, and its altitude corresponds closely to the observer's latitude. Consequently, as one moves northward, the angle of Polaris above the horizon rises, reflecting the increasing latitude. This relationship is a direct consequence of the spherical nature of the Earth and the fixed position of Polaris in the celestial sphere, which serves as a reliable navigational reference for determining latitude. Thus, the higher position of Polaris in the sky at northern latitudes exemplifies fundamental astronomical principles related to celestial navigation and the geometry of the Earth.
\\ \textit{Note:} correct because the altitude of Polaris directly correlates with the observer's latitude, resulting in an increase in Polaris's apparent altitude as one travels northward due to the Earth's spherical shape and the star's position above the North Celestial Pole.
\item \textbf{Answer:} The loss of Mars's magnetic field significantly contributed to climatic changes on the planet by facilitating atmospheric loss through the interaction with solar wind. Without a robust magnetic field to shield it, Mars was unable to protect its atmosphere from being stripped away by charged particles emanating from the Sun. This atmospheric depletion reduced the greenhouse effect, which is essential for retaining heat, leading to a dramatic cooling of the planet's surface. As a result, conditions on Mars became inhospitable, transforming it from a potentially warm and wet environment into the cold, arid landscape we observe today. The interplay between the loss of the magnetic field and atmospheric erosion underscores the critical role that magnetic protection plays in sustaining planetary climates.
\\ \textit{Note:} correct because it accurately links the absence of Mars's magnetic field to increased atmospheric loss due to solar wind, which in turn diminished the greenhouse effect and caused significant cooling of the planet's surface.
\end{enumerate}

\vfill
\noindent\textit{End of Answer Key}
\end{document}